\documentclass[12pt]{article}

\usepackage{amsmath,amssymb,hyperref,booktabs}
\topmargin -.5cm
\textheight 21cm
\oddsidemargin -.125cm
\textwidth 17cm

\def\be{\begin{equation}}
\def\ee{\end{equation}}
\def\ie{\begin{equation}\begin{aligned}}
\def\fe{\end{aligned}\end{equation}}
\newcommand{\D}{\Delta}
\newcommand{\Tr}{\mathrm{Tr}}
\newcommand{\la}{\langle}
\newcommand{\ra}{\rangle}

\begin{document}

\begin{center}
{\Large\bf A Sharpened Konishi Contradiction}
\end{center}

\section{Purpose}

This note isolates a particularly sharp version of the puzzle described in
\texttt{notes/contradiction.md}.  The point is to separate clearly:
\begin{itemize}
\item assumptions coming from the \texttt{2507.12921} background solution,
\item assumptions coming from the representation theory of the first massive string level,
\item the concrete conclusion that then appears to contradict the Konishi strong-coupling formula.
\end{itemize}

\section{The Setup}

For a first-massive Konishi state, the expected strong-coupling expansion is
\be
\D = 2R + (\D_0-4) + \frac{2}{R} + O(R^{-2}).
\ee
So for a state with $\D_0=6$ one expects
\be
\D_{\D_0=6}=2R+2+\frac{2}{R}+O(R^{-2}).
\ee
The constant $+2$ is the point at issue.

On the background side, the first-order RR solution of \texttt{2507.12921} is
\be
F^{(1)}_{\alpha\beta}=\frac{i}{R}(\gamma^{01234})_{\alpha\beta},
\ee
equivalently the constant self-dual five-form
\be
F^{(1)}_5\propto \omega_{AdS_5}+\omega_{S^5}.
\ee
This is a singlet under the manifest finite symmetry group,
\be
F^{(1)}\in \la 0,0\ra_{AdS}\otimes \la 0,0\ra_{S}.
\ee

\section{Assumptions}

The contradiction uses the following four assumptions.

\paragraph{A1.} The first-order background is a manifest
$SO(1,4)\times SO(5)$ singlet.

\paragraph{A2.} The bare Hamiltonian receives no odd-order correction.
In particular,
\be
H^{(1)}=0,
\ee
because the Hamiltonian sits in the NSNS sector while the background at first order is RR.

\paragraph{A3.} Therefore the first-order \emph{effective} Hamiltonian on fluctuations is built only
from the singlet RR background inserted against the zeroth-order Hamiltonian:
\be
\widehat H^{(1)}\varphi=[\psi^{(1)}\otimes H^{(0)}\otimes\varphi]'.
\ee
Since $\psi^{(1)}$ is RR while both $H^{(0)}$ and the chosen level-1 seed are NSNS,
this operator is sector-odd:
\be
\widehat H^{(1)}:\quad \text{NSNS}\leftrightarrow \text{RR}.
\ee
It preserves the manifest $SO(1,4)\times SO(5)$ quantum numbers, but it does \emph{not}
act diagonally within a pure NSNS subsector.

\paragraph{A4.} If a manifest irrep occurs with multiplicity one in the full first-massive flat-space
spectrum, and that unique state lies in the NSNS sector, then a sector-odd first-order operator preserving
the manifest symmetry has no RR state available to mix it with. Hence the entire first-order block vanishes
on that irrep.

\section{A Concrete Example}

The flat-space/Konishi comparison already produced a concrete candidate:
\be
\la 0,2\ra_{(2,0)},
\ee
that is, spin $(2,0)$ with sphere irrep $\la 0,2\ra$.

This label is special for two independent reasons.

\paragraph{Branched Konishi side.}
From the full branched Konishi table,
\be
(2,0)\times \la 0,2\ra
\ee
appears exactly once, and only at
\be
\D_0=6.
\ee

\paragraph{Flat-space side.}
From the complete first-massive flat-space decomposition,
\be
(2,0)\times \la 0,2\ra
\ee
also appears exactly once, entirely in the NSNS sector.

So within the full level-1 string spectrum there is exactly one state with these manifest quantum
numbers, and the Konishi table identifies it as a $\D_0=6$ state.

\section{The Contradiction}

Under assumptions A1--A4, the state
\be
\la 0,2\ra_{(2,0)}
\ee
cannot mix at first order:
\begin{itemize}
\item the first-order operator $\widehat H^{(1)}$ preserves manifest
$SO(1,4)\times SO(5)$ quantum numbers;
\item it is sector-odd, so an NSNS state can only mix with an RR state;
\item but there is no RR state with the same manifest label.
\end{itemize}

So the entire first-order block on this irrep vanishes.
There is no diagonal NSNS matrix element to compute, and there is no opposite-sector partner
state to generate an off-diagonal first-order splitting.

The tension is that the explicit first-order computations done so far for similarly simple NSNS seeds
tend to vanish by kinematics, especially for highest-weight null-polarized representatives.
Here the contradiction is even sharper: by sector and representation theory alone one gets
\be
\text{uniqueness at fixed manifest quantum numbers}
\quad + \quad
\text{sector-odd singlet first-order source}
\quad \Longrightarrow \quad
\widehat H^{(1)}=0 \text{ on this state},
\ee
so there is no $O(1)$ correction, whereas Konishi requires
\be
b_0=2.
\ee

\section{Why This Is Sharper Than the Original Puzzle}

The original $\D_0=6$ puzzle could always be softened by saying that perhaps the chosen seed was too
special, or perhaps the true Konishi state mixes with other states of the same manifest quantum numbers.

For the example
\be
\la 0,2\ra_{(2,0)},
\ee
that second escape route is absent at mass level~1:
the complete flat-space scan finds only one such state.
So if the first-order matrix element vanishes here as well, the problem is not a basis-choice issue
inside a degenerate subspace. It becomes a genuine contradiction between:
\begin{itemize}
\item the Konishi formula $b_0=\D_0-4$,
\item the singlet nature of the first-order RR background,
\item and the multiplicity-one representation content of the first-massive string spectrum.
\end{itemize}

\section{Possible Ways Out}

At present the contradiction can be avoided only if at least one assumption above fails.
The main logical options are:
\begin{itemize}
\item The relevant first-order effective Hamiltonian is not exhausted by the manifestly singlet operator
$[\psi^{(1)}\otimes H^{(0)}\otimes\varphi]'$.
\item The state identified by the manifest label $\la 0,2\ra_{(2,0)}$ is not actually the physical
Konishi eigenstate even though it is the only level-1 state with those manifest quantum numbers.
\item The first-order effective Hamiltonian does not preserve the naive NSNS/RR grading in the way assumed
above.
\item The dictionary between the strong-coupling coefficient $b_0$ and the SFT Hamiltonian eigenvalue needs
to be refined.
\end{itemize}

But absent one of these escape routes, the contradiction is sharp.

\section{Concrete Next Check}

The cleanest next calculation is not a diagonal matrix element: by sector grading that is absent from the
start.  Instead one should do one of the following.
\begin{itemize}
\item Prove directly that the manifest irrep $\la 0,2\ra_{(2,0)}$ has no RR representative at mass level~1.
\item Compute the full first-order operator on an explicit NSNS representative of
$\la 0,2\ra_{(2,0)}$ and verify that the result vanishes identically because there is no RR target state with
the same manifest quantum numbers.
\item Identify the precise assumption above that fails if the Konishi value $b_0=2$ is nevertheless to be
recovered.
\end{itemize}

If none of these steps reveals an escape route, the contradiction is immediate.

\end{document}
